\documentclass[11pt]{article}
\usepackage{icmlko}
\begin{document}

\title{스물다섯 개가 움직이고 있었다 --- 그리고 내 규칙은 그중 하나만 맞혔다}
\author{월드모델 실험실}
\date{노트 494 · 2026년 8월 3일}
\maketitle

\begin{abstract}
노트 493 이 애니에서 ``덮음이 뛰면 그 칸에서 축이 약해진다''를 두 사례로
말했다. 전 도메인 $\cdot$ 전 축에서 세 봤다.

\textbf{덮음이 학습과 유보 사이에 $15\%$p 넘게 움직이는 (도메인, 축) 짝이
$\mathbf{25}$ 개다} --- 세계애니 $7$ · 애니 $4$ · 웹툰 $4$ · 모바일 $4$
· 도서 $3$ · 아이돌 $2$ · 만화 $1$. 두 해 동안 아무도 안 세어 봤다.

\textbf{그리고 내 규칙은 반증됐다.} 덮음 차와 $|rho|$ 차의 스피어만이
$\mathbf{+0.092}$ 로 관계가 없다. 방향은 오히려 반대다 --- 덮음이 오른
$19$ 짝은 $|rho|$ 가 평균 $+0.0556$ \emph{오르고}, 내린 $5$ 짝이
$-0.0901$ 로 나쁘다.

\textbf{애니에서 잰 것은 그대로 선다} --- 거기서는 제목 길이로 같은 덮음을
맞춰 짝을 지었고, 늘 붙던 행은 $0.1914 \to 0.1773$ 으로 안 늙었다. 틀린
것은 \emph{측정}이 아니라 \emph{일반화}다. 두 사례로 규칙을 만들었다.
\end{abstract}

\begin{figure}[t]
\centering
\includegraphics[width=\textwidth]{figs/f1.pdf}
\caption{\textbf{가} --- 덮음 이동과 축 세기 변화. \textbf{나} --- 제일
많이 무너진 여덟.}
\end{figure}

\section{애니보다 큰 게 둘 있다}

무너짐 순위를 매기니 애니 트렌드는 $3$ 등이다.

\begin{center}
\begin{tabular}{llrrr}
\hline
도메인 & 축 & 덮음 차 & 학습 $|rho|$ & 유보 $|rho|$ \\
\hline
아이돌 & \texttt{target\_breadth} & $-29.2$\%p & $0.3703$ & $\mathbf{0.0428}$ \\
만화 & \texttt{goods\_scale} & $-58.7$\%p & $0.3932$ & $\mathbf{0.1107}$ \\
애니 & \texttt{trend\_level} & $+40.1$\%p & $0.1019$ & $0.0133$ \\
세계애니 & \texttt{goods\_scale} & $-20.6$\%p & $0.2156$ & $0.0927$ \\
\hline
\end{tabular}
\end{center}

앞의 둘은 \textbf{속성 표지}이고 \textbf{덮음이 내려간} 쪽이다. 노트 489
가 속성 표지에 ``격차 몫 $27\%$ 아래''를 매겼는데, 그건 \emph{붙은 행과
안 붙은 행의 차}에 대한 이야기였고 여기는 \emph{붙은 행 안에서 축이 얼마나
듣나}다. \textbf{다른 물음이고 답도 다르다.}

\section{가드는 방향을 주장하지 않는다}

\texttt{lab/loop.py} 에 \texttt{\_coverage\_shift} 를 배선했다(cfg
``덮음 이동''). 이동이 있는 짝을 \textbf{찍기만} 한다 --- 그 짝에서 학습
$|rho|$ 와 유보 $|rho|$ 를 견주려면 \emph{먼저 같은 덮음으로 짝을 맞춰야
한다}는 표시다.

\begin{quote}
\textbf{조항} --- 두 사례로 규칙을 만들었으면 \textbf{전수로 세기 전까지
규칙이라 부르지 않는다.} 노트 133 의 사촌 --- 거기서는 첫 양수를
의심했고, 여기서는 \emph{두 번째 같은 방향}을 의심한다. 두 번째가 첫
번째보다 위험하다. 하나는 우연으로 보이지만 둘은 규칙으로 보인다.
\end{quote}

\section*{남은 것}

\begin{itemize}
\item \textbf{아이돌 \texttt{target\_breadth} $0.3703 \to 0.0428$} ---
$25$ 짝 중 제일 크다. 덮음이 $95.9\% \to 66.7\%$ 로 내려갔다.
\item \textbf{만화 \texttt{goods\_scale} $0.3932 \to 0.1107$} --- 덮음
$80.0\% \to 21.3\%$. 만화는 판 무게가 작지 않다.
\item \textbf{덮음이 \emph{내려가는} 쪽이 더 나쁜가} --- $5$ 짝뿐이라
못 말한다. 세는 문턱을 $10\%$p 로 낮추면 짝이 는다.
\item \textbf{애니 덮음이 왜 $2025$ 에 뛰었나} --- 노트 493 부터 두
노트째 못 댔다.
\end{itemize}

\end{document}
